\section{Proof status, audit findings, and reproducible verification}
\label{sec:audit}

\subsection{What has been proved, and what is inherited}
The distinction between a proved identity, a finite representation,
and a numerical candidate is essential to this continuation. The
following table records the status of its principal components.

\begin{longtable}{@{}>{\raggedright\arraybackslash}p{.26\linewidth}
>{\raggedright\arraybackslash}p{.19\linewidth}
>{\raggedright\arraybackslash}p{.49\linewidth}@{}}
\toprule
Result & Status here & Precise scope\\
\midrule
\endhead
Dougall summation and the zero-excess binomial constant &
Classical input &
Dougall supplies \eqref{dg:sum}; the constant
\eqref{dg:quarticN0} is already in the Ramanujan--B\"uhring literature.\\
\addlinespace
Centered coefficients and every integer resonance &
Proved continuation &
Theorems~\ref{dg:master}, \ref{dg:closed}, and \ref{dg:alljets}
hold in the stated positive parameter domain, including residue zeros.\\
\addlinespace
Harmonic moment identities and primitives &
Proved specializations &
All displayed sums are single convergent brackets. The primitive is
coefficientwise regular at moving lower Gamma zeros.\\
\addlinespace
Nonlinear Stieltjes coordinate change &
Proved continuation &
Theorems~\ref{nc:universal}, \ref{nc:composition},
\ref{nc:Stieltjes}, and \ref{nc:rigidity} supply all local coefficients
and their sharp vanishing range.\\
\addlinespace
Three unequal integer frequencies &
Proved continuation &
Theorem~\ref{tr:finite-transport} applies to disjoint singular grids.
Its finite reduction includes all prescribed coordinate terms.\\
\addlinespace
Weighted colored kernels and Gamma correlations &
Proved consequences &
Finite root filtering and \eqref{tc:Gamma-formula} give explicit
identities; no minimal-basis or independence claim follows.\\
\addlinespace
Current Gaussian $S_6$ and $S_8$ evaluations &
Remain conjectural &
No exact functional certificate for either short evaluation is
provided here. Their inherited numerical proximity is not equality.\\
\bottomrule
\end{longtable}

The first three source questions identified in Section~\ref{sec:scope}
are answered under explicit hypotheses. The broader collision problem,
minimal arithmetic bases for spectral derivatives, and the Gaussian
evaluation conjectures remain separate questions. The current manuscript
already treats $S_4$ as proved; it must not be reset to conjectural status
when material from older arrivals is integrated \cite{repo}.

\subsection{A verified editorial correction already proposed upstream}
The canonical integration chapter contains the wording
\emph{Basis atoms must be $\Q$-independent} in its discussion of PSLQ
\cite[chapters/07-integration.tex]{repo}. The incoming Gauss--Hurwitz
report already objects to that wording \cite[Section 8]{incoming-gh}.
That correction should be retained, with credit to the earlier report.
An appropriate replacement is:
\begin{quote}
Known rational dependencies among the supplied constants should be
removed or modeled explicitly. Otherwise PSLQ may return a pre-existing
relation with zero target coefficient instead of an evaluation of the
target.
\end{quote}
This expresses the practical requirement without imposing a prior proof
of rational independence on a discovery algorithm. A relation among
the supplied atoms is a valid output of an integer-relation procedure
\cite{pslq};
it simply may not evaluate the additional target. The change is
editorial and logical, not a new counterexample to a theorem.

\subsection{Normalization points exposed by the new proofs}
No new false theorem was found in the particular incoming results used
as premises. The following are precise implementation and integration
requirements established by this article.

\begin{enumerate}
\item \textbf{Differentiate the deformation before deleting zeros.}
Equation~\eqref{dg:revivalpi} gives an exact diagnostic:
$r_0(-1)=0$ but its first spectral coefficient is $\pi^2/8$.
The regularized Gamma product or the finite polynomial shift must be
formed before taking a resonant derivative.

\item \textbf{Specify the sum and its reference coordinate.}
The subtraction in \eqref{dg:RK} uses $n+a/2$. Replacing this reference
without the finite transport \eqref{dg:reference} changes the
counterterms. In the correlation problem, the prescribed ambient
finite part gives the nonzero scalar term in
\eqref{tr:gamma-lift}; it cannot be replaced by a
zero-mean pullback without its coordinate correction.

\item \textbf{Distinguish a pointwise derivative from a distributional
derivative.} The nonlinear coefficient for
$T_{n,r}=\FP\gamma_n^{(r)}$ is $\alpha_{n,r}$, whereas that for
$D^rG_n$ is $b_{n,r}$. Their difference is the exact contact term
\eqref{nc:derivative-contact}. The derivative comparison uses the
common smooth completion of $\gamma_n(1+x)$ across the endpoint.
Extending that entire smooth remainder by zero would introduce
additional step-function derivatives and define a different object.

\item \textbf{Keep the transformed germ in composition.}
The last term of \eqref{nc:cocycle} is
$\mathfrak A_f[h\circ g]$. Leaving $h$ unchanged there would
discard the logarithmic and principal-part coefficients generated
by the first coordinate change.

\item \textbf{Separate a negative-order value from its transverse
derivative.} Equation~\eqref{tc:weighted-negative} may be
differentiated in $A$ and $B$. A derivative in $C$ needs additional
information, supplied for the ordinary kernel by
\eqref{tc:normal-derivative}.

\item \textbf{Do not silently merge singular grids.}
The full unequal-frequency theorem assumes disjoint singular sets.
At collisions, products have higher pole order and the spectral
integral has a different convergence domain. The coordinate law for
a specified coincident product in Corollary~\ref{nc:coincident}
does not by itself identify the collision limit of separated
finite parts.
\end{enumerate}

These requirements do not identify hypothetical source errors. They
state which objects the present proofs evaluate and which tempting
changes would evaluate something else.

\subsection{Finite verification and its evidentiary role}
The supplied Python programs use exact SymPy algebra and mpmath
arithmetic. The recorded runs use SymPy~1.14.0 and mpmath~1.3.0.
Their machine-readable results are supplied with the sources.

\paragraph{Dougall and harmonic branch.}
\src{code/verify_dougall.py} checks the coefficient differential
law and residue polynomial exactly through coefficient index two,
as well as the explicit quartic $C_1$. Its 19 numerical comparisons
include generic resonances through $N=3$, first and second resonant
jets, a zero-residue case, the revived finite head, quartic moments,
squared-binomial moments and harmonic sums, and
\eqref{dg:quarteridentity}. The run uses 75 decimal digits, the first
120 summands, and a finite centered asymptotic tail through $C_{11}$.
The largest observed absolute discrepancy is
$3.76169\times10^{-36}$, below the configured $10^{-31}$ threshold.
The asymptotic-tail calculation is a diagnostic approximation; it
does not furnish a rigorous remainder enclosure.

A separate unaccelerated check of \eqref{dg:quarticN1jet} gives
partial-sum errors of about $9.41\times10^{-6}$ at 100 terms and
$2.03\times10^{-10}$ at 30,000 terms against the right side
\[
-0.0147807743355189137009913877658542374\ldots.
\]
This checks the summand directly without reusing the tail completion.
The exact proof remains Theorem~\ref{dg:alljets}.

\paragraph{Nonlinear coordinate branch.}
\src{code/verify_nonlinear.py} records 138 exact finite-algebra
checks of the coefficient generator, cocycle polynomial,
low-order formulas, and filtration. The independent implementation
\src{code/independent_cutoff.py} constructs the inverse coordinate
series, transforms the test function with its inverse Jacobian,
integrates each singular monomial in the source coordinate, and
extracts the constant after composing the primitive. Its actual-cutoff
side does not use the proposed residue formula or the
$\alpha_{n,r}$ generator.

For
\[
 f(t)=q\left(t+\frac23t^2-\frac15t^3+\frac17t^4\right),
 \qquad q=1,2,
\]
it checks 48 pairs of Stieltjes and derivative indices, through
$r=3$ and $n=2r+2$, against 140 monomial test pairings.
The $q=2$ comparisons retain exact polynomial expressions in
$\log2$. Every pairing also agrees with the reusable
coordinate-coefficient implementation. These tests independently
verify signs, Jacobians, and scale terms in low degrees.

\paragraph{Unequal-frequency branch.}
\src{code/verify_transport.py} records four exact polynomial
integrals and 13 numerical comparisons at 45 decimal digits.
The numerical checks include pointwise Stieltjes multiplication,
the nonzero ambient means, direct finite-part quadrature at
$(p_1,p_2,p_3)=(1,2,3)$, permutation symmetry, the simplified
twelfth-root lower term, and the six-region Fourier normalization.
For the latter, at orders $(2,2,2)$ the Hurwitz factors are
Bernoulli polynomials. Direct piecewise polynomial integration
for the recorded shifts and frequencies gives exactly
\[
 -\frac{1879}{125411328}.
\]
A separate Mellin evaluation of the colored kernels agrees with it.
The largest recorded scaled numerical discrepancy is about
$1.46\times10^{-45}$.

\paragraph{A numerical differentiation issue.}
During verification, generic numerical differencing of
$\zeta(s,1/4)$ at $s=0$ for the second order derivative lost
accuracy in the selected backend. Native Hurwitz spectral
derivatives, corroborated by Cauchy-circle differentiation, agreed
with the independent harmonic sum. The supplied program therefore
uses the native spectral derivative interface. This is a
computational observation from the recorded run, not a correction
to the analytic identity or a claim about every implementation.

The exact all-index results are proved by finite asymptotic
remainders, normal convergence, analytic coefficient extraction,
residues, and finite root filtering. Finite symbolic tests and
floating-point agreement provide separate implementation evidence.
No numerical file in this package is labeled an interval certificate.

\subsection{Proposed repository integration}
The modular TeX sources use the label prefixes \src{dg:},
\src{dp:}, \src{nc:}, \src{tr:}, and \src{tc:}. The source manifest
pins the inspected baseline and records the incoming archive hashes.
The companion \src{INTEGRATION_NOTES.md} gives a source-to-result
crosswalk and proposed insertion points.

A suitable integration order is: retain the incoming Gauss theorem,
append the centered Dougall extension and harmonic jets; place the
primitive beside the existing Lerch and integration material; extend
the affine finite-part section with the nonlinear law; and append the
three-frequency transport after the existing pair and translated
triple formulas. Each insertion should preserve its explicit
subtraction convention and proof status. The source reports remain
provenance, while the shared lemmas need appear only once in the
consolidated manuscript.
